\section{Introducción}
El desarrollo de software actual se enfrenta a una rápida evolución tecnológica y a proyectos cada vez más complejos, lo que exige soluciones que garanticen escalabilidad, calidad y una colaboración eficaz entre equipos distribuidos. Estos retos influyen directamente en la estabilidad del código y en su capacidad de adaptación a lo largo del tiempo. Ante este panorama, la comunidad profesional y académica ha promovido el uso de herramientas y metodologías que faciliten la organización del trabajo y la trazabilidad de los cambios.

Dentro de este marco, el sistema de control de versiones Git y la plataforma colaborativa GitHub se han consolidado como pilares para la gestión de código, permitiendo registrar, revertir y coordinar modificaciones de forma segura. Paralelamente, la Guía de Scrum 2020 refuerza el enfoque ágil mediante ciclos de trabajo cortos e iterativos que impulsan la transparencia, la mejora continua y la entrega de valor en cada fase del proyecto. En el ámbito de la arquitectura de software, los patrones de diseño ofrecen soluciones probadas a problemas recurrentes, fomentando la modularidad y la reutilización del código, mientras que el Modelo-Vista-Controlador (MVC) establece una clara separación de responsabilidades que facilita el mantenimiento y la escalabilidad de las aplicaciones. Finalmente, el lenguaje Java, con su orientación a objetos y su portabilidad, se mantiene como una de las opciones más robustas y versátiles para el desarrollo de aplicaciones web, móviles y empresariales. En este contexto, el presente trabajo analiza y articula estos cinco elementos clave con el propósito de demostrar que su integración fortalece la calidad, la colaboración y la capacidad de respuesta en proyectos de software, ofreciendo un modelo de buenas prácticas aplicable en entornos académicos y profesionales.